% =======================================================================================
\section{Optimisation, Schedules and Compute}
\label{app:training}
% =======================================================================================

\subsection{Optimiser}

Gradients are clipped to a global norm of 1, then the two-dimensional weight matrices are
updated with Muon \citep{jordan2024muon,liu2025muon} and all other parameters with Adam, both
at learning rate $\gamma = 10^{-4}$ and without weight decay (\texttt{optax.contrib.muon},
version 0.2.8, default width scaling). Batches hold 256 windows; one epoch is 1000 steps.
Muon orthogonalises each matrix update, so its size per element is fixed by the learning rate
and the shape: we measured an update RMS of $(0.03$--$0.045)\,\gamma$, about five times smaller
than AdamW at the same nominal rate. We did not tune the learning rate; with the compute
available, every run went into a different question (\cref{app:xs}).

\subsection{Schedules}
\label{app:training-schedules}

\paragraph{Auxiliary heads.}
The weight $a(\tau)$ of \cref{eq:loss} decays as a half cosine from 1 to 0 over the first
$f_{\rm aux}$ of training: half of the run in the original configuration, a tenth in the
recipe.

\paragraph{Learning rate.}
The original runs kept the learning rate constant. The recipe adds a cooldown: constant, then
a linear decay to zero over the last 20\% of the steps, the warmup-stable-decay schedule of
\citet{hu2024minicpm} and \citet{hagele2024scaling} without warmup. It is applied as a scale
on the update in epoch $e$ of $E$,
\begin{equation}
  \lambda_e = \begin{cases} 1 & e < E - E_c,\\[2pt] \dfrac{E - e - \frac12}{E_c} & e\ge E-E_c,\end{cases}
\end{equation}
with $E_c$ the number of cooldown epochs, so that each epoch takes the ramp's value at its
midpoint. Without weight decay, scaling the update is exactly scaling the learning rate, and
because the scale is a function of the epoch rather than of a step counter in the optimiser
state, a run resumed with a fresh optimiser lands at the right point of the schedule. The
cooldowns of \cref{app:xs} continue a finished 1M-step run for 200k more steps with
$E_c = 200$.

\paragraph{Checkpoints.}
Every epoch saves a checkpoint, and a job longer than the cluster's 24-hour limit is a chain
of jobs that each resume from the latest one. To fit a storage quota, the 768-wide runs and
B save the weights without the optimiser state. A resumed job then restarts Muon's momentum
and Adam's moments, which leaves no visible trace in the loss: on B, the epoch before and
after the first restart logged 0.5370 and 0.5374.

\subsection{Compute}
\label{app:training-compute}

Every run used a single NVIDIA A100 (80\,GB, PCIe). The simulator, the loss and the update
are compiled into one XLA program per epoch of 1000 steps. \Cref{tab:compute} gives the
measured costs. A fit of the step time against the network's FLOPs on XS and S gives about
23\,ms of fixed cost per step and 105\,TFLOP/s of marginal throughput (about a third of the
A100's bfloat16 peak), so the small XS problem is dominated by overheads and the large B
problem by the network. The whole XS ablation ladder of \cref{tab:xs-ladder} cost about 45
A100-hours, and B1 about 86.

\begin{table}[h]
\caption{Measured cost on one A100, batch 256. Peak memory is for training; JAX reserves 75\%
of the card by default.}
\label{tab:compute}
\centering
\small
\begin{tabular}{llrrl}
\toprule
rung & width & per 1000 steps & peak memory & runs \\
\midrule
XS & 512 & 38.9\,s & -- & 1M steps in 10.8\,h; 200k cooldown in 2.2\,h \\
XS & 768 & 59.6\,s & -- & 1M steps in 16.6\,h; 200k cooldown in 3.6\,h \\
S & 512 & 80\,s & -- & 500k steps in 11.1\,h (an early run that diverged; not retrained) \\
B & 512 & 1536\,s & 18.6\,GiB & 200k steps in 86\,h (4 chained jobs) \\
B & 768 & 2570\,s & 29.4\,GiB & 300-step benchmark \\
\bottomrule
\end{tabular}
\end{table}

\paragraph{Sampling cost.}
A scorecard pass on XS (210 windows, 256 draws each, 32 RK4 steps) takes about 4 minutes on
an A100, about one second per window. On B, 256 draws of one window take about 78 seconds with
the 768-wide network.

\subsection{Numerics}
\label{app:training-numerics}

Training and evaluation run in bfloat16. Three checks, on XS-late
(\cref{tab:numerics}), show that the numerics do not set the results: the RK4 integrator is
converged at 32 steps; float32 evaluation on a CPU gives the same widths as bfloat16 on the
A100 (median ratio 0.991, 95\% interval 0.978--1.000, over 93 loud sources) and finds the same
169 sources; and float32 on the A100 reproduces the CPU once XLA's GEMM autotuning is switched
off. With autotuning on, float32 programs were silently wrong on both GPUs we used. At the
default level they found almost no sources. At level 3 the A100 gave posteriors 2.4 times
narrower than the CPU's at the same distance from the truth: over-confident, with 45\% of the
truths inside the 68\% intervals. A fixed compiled program gave different answers at different
autotuning levels, on two JAX and two CUDA versions, so the fault is in the compilation and not
in the model; bfloat16 is unaffected. We report it because the over-confident failure is easy to
mistake for a sharper posterior.
